\chapter{\texorpdfstring{\(L^2\)}{L2} Fourier 级数}
\label{chap:l2-fourier-series}

点态收敛观察每个 \(x\) 上的数值，\(L^2\) 收敛观察一个周期上的平均平方
误差。后一问题与 Hilbert 空间的正交投影完全一致：三角部分和是到有限维
三角多项式空间的投影，唯一尚需证明的是三角系统没有遗漏方向。Fejér 定理
给出的三角多项式稠密性正好完成这一步。

本章沿用
\[
 \langle f,g\rangle_{\mathbb T}
 =\frac1{2\pi}\int_{-\pi}^{\pi}f(x)\overline{g(x)}\,dx
\]
以及\chapterref{chap:trigonometric-fourier-series}的 Fourier 系数约定。

\section{三角系统的完备性}

函数 \(e_k(x)=e^{ikx}\) 满足
\(\langle e_j,e_k\rangle_{\mathbb T}=\delta_{jk}\)，所以
\((e_k)_{k\in\mathbb Z}\) 是 \(L^2(\mathbb T)\) 中的正交规范系。

\begin{theorem}\label{thm:trigonometric-system-complete}
三角系统 \((e^{ikx})_{k\in\mathbb Z}\) 在 \(L^2(\mathbb T)\) 中完备。
\end{theorem}
\begin{proof}
设 \(f\in L^2(\mathbb T)\) 与每个 \(e_k\) 正交，则全部 Fourier 系数为零。
由于圆周测度有限，\(f\in L^1(\mathbb T)\)。
\chapterref{chap:trigonometric-fourier-series}的 Fejér 平均
\(\sigma_Nf\) 的系数是
\[
 \widehat{\sigma_Nf}(k)=
 \begin{cases}
 (1-|k|/(N+1))\widehat f(k),&|k|\le N,\\
 0,&|k|>N,
 \end{cases}
\]
故 \(\sigma_Nf=0\)。另一方面，Fejér 核的 \(L^2\) 逼近恒等性质给
\(\sigma_Nf\to f\) 于 \(L^2\)：这可直接调用
\chapterref{chap:trigonometric-fourier-series}的 \(L^p\) 结论，取 \(p=2\)。
因此 \(f=0\)。由正交系完备判据，三角系统完备。
\end{proof}

这里没有从 Fourier 系数唯一性循环推出完备性。唯一性在 \(L^1\) 中由 Fejér
平均证明，而 Fejér 收敛只使用正核和连续函数稠密性，不预设三角系统完备。

\section{\texorpdfstring{\(L^2\)}{L2} 收敛}

\begin{theorem}[Riesz--Fischer 定理的 Fourier 形式]
\label{thm:l2-fourier-convergence}
对每个 \(f\in L^2(\mathbb T)\)，
\[
 S_Nf=\sum_{k=-N}^{N}\widehat f(k)e^{ikx}
 \longrightarrow f\qquad\text{于 }L^2(\mathbb T).
\]
反之，对每个 \(c=(c_k)_{k\in\mathbb Z}\in\ell^2(\mathbb Z)\)，存在唯一
\(f\in L^2(\mathbb T)\) 使 \(\widehat f(k)=c_k\) 对所有 \(k\) 成立。
\end{theorem}
\begin{proof}
第一部分是完备正交系展开定理。对第二部分，正交级数
\(\sum_kc_ke^{ikx}\) 在 \(L^2\) 中收敛到某个 \(f\)。与固定 \(e_j\) 取内积
是连续操作，故极限的第 \(j\) 个系数为 \(c_j\)。若另一个函数有相同系数，
两者之差与完备系统正交，因而在 \(L^2\) 中为零。
\end{proof}

这个结论不声称 \(S_Nf(x)\) 在每一点收敛。\(L^2\) 元素本来就是几乎处处
等价类，改变零测集上的值不会影响结论；而点态收敛会看见这些代表元差别。
对一般 \(L^2\) 函数，点态收敛需要比本章多得多的工具。

\section{Parseval 恒等式}

\begin{corollary}[Parseval 恒等式]\label{thm:fourier-parseval-circle}
对 \(f,g\in L^2(\mathbb T)\)，
\[
 \|f\|_{L^2(\mathbb T)}^2=\sum_{k\in\mathbb Z}|\widehat f(k)|^2,
 \qquad
 \langle f,g\rangle_{\mathbb T}
 =\sum_{k\in\mathbb Z}\widehat f(k)\overline{\widehat g(k)}.
\]
\end{corollary}

因此 Fourier 变换
\[
 \mathcal F_{\mathbb T}:L^2(\mathbb T)\to\ell^2(\mathbb Z),
 \qquad f\mapsto(\widehat f(k))_{k\in\mathbb Z}
\]
是酉同构。部分和误差有精确表达式
\begin{equation}\label{eq:l2-fourier-tail}
 \|f-S_Nf\|_2^2=\sum_{|k|>N}|\widehat f(k)|^2.
\end{equation}
这比只知道误差趋零更强：任何系数衰减估计都会立即转化为逼近速度。

\section{平移、微分与系数衰减}

对 \(h\in\R\)，平移 \(\tau_hf(x)=f(x-h)\) 满足
\[
 \widehat{\tau_hf}(k)=e^{-ikh}\widehat f(k),\qquad
 \|\tau_hf-f\|_2^2
 =\sum_k|e^{-ikh}-1|^2|\widehat f(k)|^2.
\]
控制收敛定理在计数测度上给出 \(L^2\) 平移连续性。

若 \(f\) 是 \(2\pi\)-周期绝对连续函数且 \(f'\in L^2\)，分部积分给
\begin{equation}\label{eq:fourier-derivative-coefficients}
 \widehat{f'}(k)=ik\widehat f(k).
\end{equation}
边界项因周期性消失。Parseval 恒等式于是给
\[
 \|f'\|_2^2=\sum_{k\in\mathbb Z}k^2|\widehat f(k)|^2.
\]
反复应用可知：若有 \(r\) 阶平方可积导数，则
\(|k|^r\widehat f(k)\in\ell^2\)。这是一种平均意义的衰减，不必逐项推出
\(|\widehat f(k)|\le C|k|^{-r-1}\)。

反向也成立：若 \(f\in L^2\) 且
\(\sum k^2|\widehat f(k)|^2<\infty\)，则存在 \(g\in L^2\) 的系数为
\(ik\widehat f(k)\)，并且
\[
 \frac1{2\pi}\int_{-\pi}^{\pi} f\,\varphi'
 =-\frac1{2\pi}\int_{-\pi}^{\pi}g\,\varphi
\]
对所有光滑周期函数 \(\varphi\) 成立。此时 \(g\) 是 \(f\) 的弱导数。
\chapterref{chap:distributions-weak-derivatives}和
\chapterref{chap:sobolev-spaces}将把这一定义推广到开集和高维情形。

\section{卷积与乘子}

圆周卷积定义为
\[
 (f*g)(x)=\frac1{2\pi}\int_{-\pi}^{\pi}f(x-t)g(t)\,dt.
\]
若 \(f\in L^p,g\in L^1\)，Young 不等式给
\(\|f*g\|_p\le\|f\|_p\|g\|_1\)，其中范数均相对归一化测度。Fubini 定理给
\begin{equation}\label{eq:circle-convolution-coefficients}
 \widehat{f*g}(k)=\widehat f(k)\widehat g(k).
\end{equation}
所以卷积在 Fourier 坐标中成为逐项乘法。

更一般地，给定有界数列 \(m=(m_k)\)，定义
\[
 T_mf=\sum_{k\in\mathbb Z}m_k\widehat f(k)e^{ikx}.
\]
Parseval 恒等式给
\(\|T_m\|_{L^2\to L^2}=\sup_k|m_k|\)。这类算子称为 Fourier 乘子。
部分和对应 \(m_k=\mathbf1_{\{|k|\le N\}}\)，Fejér 平均对应
\(m_k=(1-|k|/(N+1))_+\)。在 \(L^2\) 中二者都是收缩；但在
\(L^\infty\) 中，Dirichlet 核的 \(L^1\) 增长使部分和没有共同算子界。

\section*{习题}
\addcontentsline{toc}{section}{习题}
\subsection*{A组}
\begin{enumerate}[label=\textbf{\thechapter.\arabic*.}]
\exerciseitem{l2-fourier-orthonormality}验证 \(e^{ikx}\) 在归一化内积下正交规范。
\exerciseitem{l2-fourier-best-approximation}证明 \(S_Nf\) 是所有次数不超过 \(N\) 的三角多项式中唯一的 \(L^2\) 最佳逼近。
\exerciseitem{l2-fourier-tail-error}由 Parseval 恒等式推导误差公式\eqref{eq:l2-fourier-tail}。
\exerciseitem{l2-fourier-translation-continuity}只用系数和控制收敛证明 \(\|\tau_hf-f\|_2\to0\)。
\end{enumerate}
\subsection*{B组}
\begin{enumerate}[label=\textbf{\thechapter.\arabic*.},resume]
\exerciseitem{l2-fourier-sawtooth-parseval}求 \(f(x)=x\)（\(-\pi<x<\pi\)）的 Fourier 系数，并由 Parseval 恒等式计算 \(\sum_{n=1}^\infty n^{-2}\)。
\exerciseitem{l2-fourier-derivative-converse}补全“\((k\widehat f(k))\in\ell^2\) 蕴含存在弱导数”的证明。
\exerciseitem{l2-fourier-poincare-circle}若 \(\widehat f(0)=0\) 且 \(f\) 有平方可积弱导数，证明 \(\|f\|_2\le\|f'\|_2\)，并讨论等号。
\exerciseitem{l2-fourier-convolution-young}证明圆周上的 Young 不等式及卷积系数公式。
\exerciseitem{l2-fourier-multiplier-norm}证明 \(L^2\) Fourier 乘子的范数恰为 \(\sup_k|m_k|\)。
\end{enumerate}
\subsection*{C组}
\begin{enumerate}[label=\textbf{\thechapter.\arabic*.},resume]
\exerciseitem{l2-fourier-compact-operator}若 \(m_k\to0\)，证明 \(T_m\) 是 \(L^2(\mathbb T)\) 上的紧算子；反之亦然。
\exerciseitem{l2-fourier-subsequence-ae}由 \(L^2\) 收敛的子列原理证明：每个 \(f\in L^2(\mathbb T)\) 都有一列三角多项式几乎处处收敛到 \(f\)。
\optionalexerciseitem{l2-fourier-autocorrelation}对 \(f\in L^2(\mathbb T)\) 定义 \(R_f(h)=\langle \tau_hf,f\rangle\)。证明 \(R_f\) 连续且正定，并求其 Fourier 系数。
\end{enumerate}
